% 第4章 监督学习与 alpha 挖掘
\section{监督学习与 alpha 挖掘}\label{sec:supervised}

监督学习是最早进入量化、也最容易被低估的一族。它做的事情听起来很朴素：给定一批特征，预测下一期的横截面收益，然后按预测排序构建多空组合。可正是这个最朴素的问题，逼出了金融机器学习里最本质的几组矛盾——低信噪比下的过拟合、预测目标与交易目标的错位、以及“因子被发表后收益衰减”。

\subsection{问题设定}

形式化地，在每个调仓时点 $t$，观察到每只资产 $i$ 的特征向量 $x_{i,t}\in\mathbb{R}^p$（估值、动量、质量、流动性等），要预测其下一期收益 $r_{i,t+1}$。学习目标是
\begin{equation}
  \hat f = \arg\min_{f\in\mathcal{H}} \sum_{t}\sum_{i} \mathcal{L}\big(r_{i,t+1},\, f(x_{i,t})\big) + \Omega(f),
\end{equation}
$\mathcal{L}$ 是损失（均方或排序），$\Omega$ 是正则项。这一步看似标准，但\textbf{目标函数与“赚钱”之间隔着三层}：预测的是收益而非风险调整收益；以横截面为单位而非以组合为单位；没有把交易成本写进去。三层之外还有一层的错位——调仓频率决定了标签的构造（未来 1 期还是 $k$ 期），而标签构造会反过来决定模型学到的是动量还是反转。

\subsection{从线性到非线性：一条可复现的演进线}

\textbf{线性基准}。Fama--French 三因子/五因子模型是这一族的天花板之一，也是所有非线性模型必须超越的基线。它们本质上是把 $\hat f$ 限制成少数暴露的线性组合。符号上简单，但“用哪些暴露、怎么组合”始终是个开放问题。

\textbf{正则化回归}。当候选特征从几个扩到几百个，问题立刻变成高维信号的挑选。Feng、Giglio 与 Xiu（2020）\cite{fenggiglioxu2020} 用双重选择 Lasso（double-selection）来处理“因子动物园”问题，思路是先控制住一大片潜在因子，再判断新因子是否有增量，这与 Harvey、Liu 与 Zhu（2016）\cite{harveylz2016} 的多重检验警示是同一件事的两面。

\textbf{树模型与集成}。梯度提升树（GBRT）与随机森林在金融横截面上的表现常好于线性模型，原因是它们能自然处理非线性和交互，并对异常值不敏感。Krauss、Do 与 Huck（2017）\cite{krauss2017} 在标普 500 成分股上系统比较了深度神经网络、梯度提升树与随机森林三种方法做统计套利，发现集成方法在多数设定下稳定占优——这是把“到底哪种模型好”放到公开口径下比较的早期代表作之一。

\textbf{神经网络}。Gu、Kelly 与 Xiu（2020）\cite{gukellyxiu2020} 是这一族的基准文献。他们用同样的特征集，把线性回归、正则化回归、GBRT、随机森林与多种神经网络放在一起比较，报告的样本外月度 $R^2$ 只有约 0.3\%--0.4\%，但由此构建的多空组合的夏普与累计收益显著。这一结果确立了这一族最重要的经验事实：\textbf{解释力极小不等于不可用，关键在于是否把微弱信号稳定地转化为组合收益。}他们还报告了神经网络相对树模型的优势，并强调这一优势来自对交互与非线性的刻画，而非模型复杂度本身。

\textbf{把无套利约束写进网络}。Chen、Pelger 与 Zhu（2024）\cite{chenpelgerzhu2024} 的做法是把资产定价的无套利结构（因子定价的线性结构）直接编码进深度网络的架构，让网络在“定价核”层面学习而非在收益层面自由拟合。这类方法的价值在于\textbf{把经济结构当作先验注入模型}，从而在低信噪比下减少需要估计的自由度——这与纯黑箱拟合是两种哲学。

\textbf{“复杂性红利”的争议}。Kelly、Malamud 与 Zhou（2024）\cite{kellymalamudzhou2024} 提出“复杂性的价值”（virtue of complexity）：在高维但结构稀疏的设定下，一个参数数量远大于样本量的岭回归反而能有更好的样本外表现。这个结论在文献里争议不小——支持者认为它挑战了“奥卡姆剃刀”，反对者指出它对信号强度与协方差结构做了较强假设，且在高频、真实的交易成本下未必成立。本文在此不站队，只把它列为\textbf{证据强度中等、边界条件明确}的一类主张。

\subsection{落地要素}

要复现这一族的任何工作，最容易被忽略、也最致命的是三件事：

\begin{itemize}
  \item \textbf{时点对齐（point-in-time）}。财报数据必须用实际披露日，而非报告期；任何“用当期财报预测当期收益”的设定都是前瞻偏差。这是这一族数据管线里最容易出错的地方。
  \item \textbf{横截面中性化}。把行业、市值、风格暴露从预测里剥掉，否则模型学到的可能只是“小盘/成长”这些已知风险溢价，而非新 alpha。
  \item \textbf{交易成本口径}。横截面策略的换手率高，成本对净值曲线的影响常常是决定性的。不报告成本的样本外夏普，与报告成本后的夏普，可能相差一个量级。
\end{itemize}

\subsection{复现状态与失效边界}

这一族是六族里复现基础最好的：Microsoft Qlib\cite{qlib} 公开了 Alpha158/Alpha360 两个特征集与一批模型在 CSI300 上的基准结果，Grinold--Kahn 风格的因子研究也有成熟的开源工具链（Alphalens 等）。但“能跑通”不等于“能赚钱”：Gu、Kelly 与 Xiu 的样本外结论在换到 A 股、换到更近的样本区间后是否成立，公开的独立复现并不多（详见第 \ref{sec:repro_audit} 章）。

失效边界主要有四条，且都可检验：

\begin{enumerate}
  \item \textbf{信息耗尽}。任何被广泛使用的特征，其边际信息会被 Kyle 型的价格冲击机制压薄（McLean \& Pontiff 2016\cite{mcleanpontiff2016} 记录了这一衰减）。
  \item \textbf{非平稳}。特征与收益的关系随市场制度漂移，固定模型会在制度切换后失效。
  \item \textbf{成本吞噬}。高换手策略在真实成本下的净收益可能为负。
  \item \textbf{数据窥探}。在几百个特征与几十个模型里挑最好的，样本内冠军往往是噪声冠军。
\end{enumerate}

本章要记住的一句话是：\textbf{监督学习在金融里的真正难点，从来不是“用什么模型”，而是“在信息耗尽、分布漂移、成本与多重检验的四重约束下，还能剩下多少稳定的、可交易的增量”。}
